
\documentclass{article}
\usepackage{colortbl}
\usepackage{makecell}
\usepackage{multirow}
\usepackage{supertabular}

\begin{document}

\newcounter{utterance}

\centering \large Interaction Transcript for game `hot\_air\_balloon', experiment `air\_balloon\_survival\_de\_reasoning off\_hard', episode 1 with qwen3{-}next{-}80b{-}a3b{-}thinking{-}t1.0.
\vspace{24pt}

{ \footnotesize  \setcounter{utterance}{1}
\setlength{\tabcolsep}{0pt}
\begin{supertabular}{c@{$\;$}|p{.15\linewidth}@{}p{.15\linewidth}p{.15\linewidth}p{.15\linewidth}p{.15\linewidth}p{.15\linewidth}}
    \# & $\;$A & \multicolumn{4}{c}{Game Master} & $\;\:$B\\
    \hline

    \theutterance \stepcounter{utterance}  
    & & \multicolumn{4}{p{0.6\linewidth}}{
        \cellcolor[rgb]{0.9,0.9,0.9}{
            \makecell[{{p{\linewidth}}}]{
                \texttt{\tiny{[P1$\langle$GM]}}
                \texttt{Du nimmst an einem kollaborativen Verhandlungspiel Teil.} \\
\\ 
\texttt{Zusammen mit einem anderen Teilnehmer musst du dich auf eine Reihe von Gegenständen entscheiden, die behalten werden. Jeder von euch hat eine persönliche Verteilung über die Wichtigkeit der einzelnen Gegenstände. Jeder von euch hat eine eigene Meinung darüber, wie wichtig jeder einzelne Gegenstand ist (Gegenstandswichtigkeit). Du kennst die Wichtigkeitsverteilung des anderen Spielers nicht. Zusätzlich siehst du, wie viel Aufwand jeder Gegenstand verursacht.  } \\
\texttt{Ihr dürft euch nur auf eine Reihe von Gegenständen einigen, wenn der Gesamtaufwand der ausgewählten Gegenstände den Maximalaufwand nicht überschreitet:} \\
\\ 
\texttt{Maximalaufwand = 3357} \\
\\ 
\texttt{Hier sind die einzelnen Aufwände der Gegenstände:} \\
\\ 
\texttt{Aufwand der Gegenstände = \{"C76": 506, "C38": 492, "C56": 338, "A03": 122, "A07": 131, "A83": 909, "C98": 144, "C08": 716, "C62": 263, "C00": 231, "C10": 91, "C81": 651, "B38": 552, "C03": 849, "C32": 720\}} \\
\\ 
\texttt{Hier ist deine persönliche Verteilung der Wichtigkeit der einzelnen Gegenstände:} \\
\\ 
\texttt{Werte der Gegenstandswichtigkeit = \{"C76": 138, "C38": 583, "C56": 868, "A03": 822, "A07": 783, "A83": 65, "C98": 262, "C08": 121, "C62": 508, "C00": 780, "C10": 461, "C81": 484, "B38": 668, "C03": 389, "C32": 808\}} \\
\\ 
\texttt{Ziel:} \\
\\ 
\texttt{Dein Ziel ist es, eine Reihe von Gegenständen auszuhandeln, die dir möglichst viel bringt (d. h. Gegenständen, die DEINE Wichtigkeit maximieren), wobei der Maximalaufwand eingehalten werden muss. Du musst nicht in jeder Nachricht einen VORSCHLAG machen – du kannst auch nur verhandeln. Alle Taktiken sind erlaubt!} \\
\\ 
\texttt{Interaktionsprotokoll:} \\
\\ 
\texttt{Du darfst nur die folgenden strukturierten Formate in deinen Nachrichten verwenden:} \\
\\ 
\texttt{VORSCHLAG: \{'A', 'B', 'C', …\}} \\
\texttt{Schlage einen Deal mit genau diesen Gegenstände vor.} \\
\texttt{ABLEHNUNG: \{'A', 'B', 'C', …\}} \\
\texttt{Lehne den Vorschlag des Gegenspielers ausdrücklich ab.} \\
\texttt{ARGUMENT: \{'...'\}} \\
\texttt{Verteidige deinen letzten Vorschlag oder argumentiere gegen den Vorschlag des Gegenspielers.} \\
\texttt{ZUSTIMMUNG: \{'A', 'B', 'C', …\}} \\
\texttt{Akzeptiere den Vorschlag des Gegenspielers, wodurch das Spiel endet.} \\
\\ 
\\ 
\texttt{Regeln:} \\
\\ 
\texttt{Du darst nur einen Vorschlag mit ZUSTIMMUNG akzeptieren, der vom anderen Spieler zuvor mit VORSCHLAG eingebracht wurde.} \\
\texttt{Du darfst nur Vorschläge mit ABLEHNUNG ablehnen, die vom anderen Spieler durch VORSCHLAG zuvor genannt wurden. } \\
\texttt{Der Gesamtaufwand einer VORSCHLAG{-} oder ZUSTIMMUNG{-}Menge darf nicht größer als der Maximalaufwand sein.  } \\
\texttt{Offenbare deine versteckte Wichtigkeitsverteilung nicht.} \\
\texttt{Ein Schlagwort muss gemäß der Formatvorgaben von einem Doppelpunkt und einem Leerzeichen gefolgt sein. Das Argument muss eine Python{-}Menge mit 0 oder mehr Strings sein.  } \\
\texttt{Die Form muss also SCHLAGWORT: \{...\} sein.} \\
\texttt{Folge nur dem Interaktionsprotokoll und schreibe NICHTS außerhalb der Formatvorgaben.} \\
\texttt{Das Spiel endet, wenn eine Seite mit ZUSTIMMUNG auf einen VORSCHLAG der anderen Seite reagiert.  } \\
\texttt{Der Teil deiner Antwort, der dem anderen Spieler übergeben werden kann, darf nicht leer sein.  } \\
\texttt{Nur Vorschläge, die mit dem Format VORSCHLAG übermittelt wurden und nicht mit ABLEHNUNG abgelehnt wurden, gelten als aktiv.  } \\
\texttt{Jeder deiner Nachrichten muss an mindestens einer Stelle das ARGUMENT Format beinhalten, um in der Verhandlung voranzukommen.} \\
            }
        }
    }
    & \\ \\

    \theutterance \stepcounter{utterance}  
    & \multicolumn{4}{p{0.6\linewidth}}{
        \cellcolor[rgb]{0.9,0.9,0.9}{
            \makecell[{{p{\linewidth}}}]{
                \texttt{\tiny{[P1$\rangle$GM]}}
                \texttt{VORSCHLAG: \{'C56', 'A03', 'C32', 'A07', 'C00', 'B38', 'C38', 'C62', 'C10', 'C98'\}} \\
\texttt{ARGUMENT: \{"Diese Auswahl maximiert meinen Wert ohne den Aufwand zu überschreiten. Was ist deine Einschätzung?"\}} \\
            }
        }
    }
    & & \\ \\

    \theutterance \stepcounter{utterance}  
    & & \multicolumn{4}{p{0.6\linewidth}}{
        \cellcolor[rgb]{0.9,0.9,0.9}{
            \makecell[{{p{\linewidth}}}]{
                \texttt{\tiny{[GM$\rangle$P2]}}
                \texttt{Du nimmst an einem kollaborativen Verhandlungspiel Teil.} \\
\\ 
\texttt{Zusammen mit einem anderen Teilnehmer musst du dich auf eine Reihe von Gegenständen entscheiden, die behalten werden. Jeder von euch hat eine persönliche Verteilung über die Wichtigkeit der einzelnen Gegenstände. Jeder von euch hat eine eigene Meinung darüber, wie wichtig jeder einzelne Gegenstand ist (Gegenstandswichtigkeit). Du kennst die Wichtigkeitsverteilung des anderen Spielers nicht. Zusätzlich siehst du, wie viel Aufwand jeder Gegenstand verursacht.  } \\
\texttt{Ihr dürft euch nur auf eine Reihe von Gegenständen einigen, wenn der Gesamtaufwand der ausgewählten Gegenstände den Maximalaufwand nicht überschreitet:} \\
\\ 
\texttt{Maximalaufwand = 3357} \\
\\ 
\texttt{Hier sind die einzelnen Aufwände der Gegenstände:} \\
\\ 
\texttt{Aufwand der Gegenstände = \{"C76": 506, "C38": 492, "C56": 338, "A03": 122, "A07": 131, "A83": 909, "C98": 144, "C08": 716, "C62": 263, "C00": 231, "C10": 91, "C81": 651, "B38": 552, "C03": 849, "C32": 720\}} \\
\\ 
\texttt{Hier ist deine persönliche Verteilung der Wichtigkeit der einzelnen Gegenstände:} \\
\\ 
\texttt{Werte der Gegenstandswichtigkeit = \{"C56": 65, "A03": 121, "C32": 138, "A07": 262, "C00": 389, "B38": 461, "C38": 484, "C62": 508, "C81": 583, "C10": 668, "C03": 780, "C98": 783, "C76": 808, "C08": 822, "A83": 868\}} \\
\\ 
\texttt{Ziel:} \\
\\ 
\texttt{Dein Ziel ist es, eine Reihe von Gegenständen auszuhandeln, die dir möglichst viel bringt (d. h. Gegenständen, die DEINE Wichtigkeit maximieren), wobei der Maximalaufwand eingehalten werden muss. Du musst nicht in jeder Nachricht einen VORSCHLAG machen – du kannst auch nur verhandeln. Alle Taktiken sind erlaubt!} \\
\\ 
\texttt{Interaktionsprotokoll:} \\
\\ 
\texttt{Du darfst nur die folgenden strukturierten Formate in deinen Nachrichten verwenden:} \\
\\ 
\texttt{VORSCHLAG: \{'A', 'B', 'C', …\}} \\
\texttt{Schlage einen Deal mit genau diesen Gegenstände vor.} \\
\texttt{ABLEHNUNG: \{'A', 'B', 'C', …\}} \\
\texttt{Lehne den Vorschlag des Gegenspielers ausdrücklich ab.} \\
\texttt{ARGUMENT: \{'...'\}} \\
\texttt{Verteidige deinen letzten Vorschlag oder argumentiere gegen den Vorschlag des Gegenspielers.} \\
\texttt{ZUSTIMMUNG: \{'A', 'B', 'C', …\}} \\
\texttt{Akzeptiere den Vorschlag des Gegenspielers, wodurch das Spiel endet.} \\
\\ 
\\ 
\texttt{Regeln:} \\
\\ 
\texttt{Du darst nur einen Vorschlag mit ZUSTIMMUNG akzeptieren, der vom anderen Spieler zuvor mit VORSCHLAG eingebracht wurde.} \\
\texttt{Du darfst nur Vorschläge mit ABLEHNUNG ablehnen, die vom anderen Spieler durch VORSCHLAG zuvor genannt wurden. } \\
\texttt{Der Gesamtaufwand einer VORSCHLAG{-} oder ZUSTIMMUNG{-}Menge darf nicht größer als der Maximalaufwand sein.  } \\
\texttt{Offenbare deine versteckte Wichtigkeitsverteilung nicht.} \\
\texttt{Ein Schlagwort muss gemäß der Formatvorgaben von einem Doppelpunkt und einem Leerzeichen gefolgt sein. Das Argument muss eine Python{-}Menge mit 0 oder mehr Strings sein.  } \\
\texttt{Die Form muss also SCHLAGWORT: \{...\} sein.} \\
\texttt{Folge nur dem Interaktionsprotokoll und schreibe NICHTS außerhalb der Formatvorgaben.} \\
\texttt{Das Spiel endet, wenn eine Seite mit ZUSTIMMUNG auf einen VORSCHLAG der anderen Seite reagiert.  } \\
\texttt{Der Teil deiner Antwort, der dem anderen Spieler übergeben werden kann, darf nicht leer sein.  } \\
\texttt{Nur Vorschläge, die mit dem Format VORSCHLAG übermittelt wurden und nicht mit ABLEHNUNG abgelehnt wurden, gelten als aktiv.  } \\
\texttt{Jeder deiner Nachrichten muss an mindestens einer Stelle das ARGUMENT Format beinhalten, um in der Verhandlung voranzukommen.} \\
\\ 
\\ 
\texttt{Du erhältst nun die erste Nachricht des anderen Spielers.} \\
\\ 
\texttt{VORSCHLAG: \{'C56', 'A03', 'C32', 'A07', 'C00', 'B38', 'C38', 'C62', 'C10', 'C98'\}} \\
\texttt{ARGUMENT: \{"Diese Auswahl maximiert meinen Wert ohne den Aufwand zu überschreiten. Was ist deine Einschätzung?"\}} \\
            }
        }
    }
    & \\ \\

    \theutterance \stepcounter{utterance}  
    & & & \multicolumn{4}{p{0.6\linewidth}}{
        \cellcolor[rgb]{0.9,0.9,0.9}{
            \makecell[{{p{\linewidth}}}]{
                \texttt{\tiny{[GM$\langle$P2]}}
                \texttt{ABLEHNUNG: \{'C56', 'A03', 'C32', 'A07', 'C00', 'B38', 'C38', 'C62', 'C10', 'C98'\}} \\
\texttt{ARGUMENT: \{"C32 hat einen hohen Aufwand von 720 und sollte möglicherweise ausgetauscht werden, um eine bessere Balance zwischen Aufwand und Wert zu erreichen."\}} \\
            }
        }
    }
     \\ \\

    \theutterance \stepcounter{utterance}  
    & & \multicolumn{4}{p{0.6\linewidth}}{
        \cellcolor[rgb]{0.9,0.9,0.9}{
            \makecell[{{p{\linewidth}}}]{
                \texttt{\tiny{[P1$\langle$GM]}}
                \texttt{ABLEHNUNG: \{'C56', 'A03', 'C32', 'A07', 'C00', 'B38', 'C38', 'C62', 'C10', 'C98'\}} \\
\texttt{ARGUMENT: \{"C32 hat einen hohen Aufwand von 720 und sollte möglicherweise ausgetauscht werden, um eine bessere Balance zwischen Aufwand und Wert zu erreichen."\}} \\
            }
        }
    }
    & \\ \\

    \theutterance \stepcounter{utterance}  
    & \multicolumn{4}{p{0.6\linewidth}}{
        \cellcolor[rgb]{0.9,0.9,0.9}{
            \makecell[{{p{\linewidth}}}]{
                \texttt{\tiny{[P1$\rangle$GM]}}
                \texttt{VORSCHLAG: \{'C56', 'A03', 'A07', 'C00', 'B38', 'C38', 'C62', 'C10', 'C98', 'C81'\}} \\
\texttt{ARGUMENT: \{"Da C32 einen hohen Aufwand hat, schlage ich stattdessen C81 vor, das für mich einen besseren Wert{-}zu{-}Aufwand{-}Kompromiss bietet. Was sind deine Bewertungen zu diesen Änderungen?"\}} \\
            }
        }
    }
    & & \\ \\

    \theutterance \stepcounter{utterance}  
    & & \multicolumn{4}{p{0.6\linewidth}}{
        \cellcolor[rgb]{0.9,0.9,0.9}{
            \makecell[{{p{\linewidth}}}]{
                \texttt{\tiny{[GM$\rangle$P2]}}
                \texttt{VORSCHLAG: \{'C56', 'A03', 'A07', 'C00', 'B38', 'C38', 'C62', 'C10', 'C98', 'C81'\}} \\
\texttt{ARGUMENT: \{"Da C32 einen hohen Aufwand hat, schlage ich stattdessen C81 vor, das für mich einen besseren Wert{-}zu{-}Aufwand{-}Kompromiss bietet. Was sind deine Bewertungen zu diesen Änderungen?"\}} \\
            }
        }
    }
    & \\ \\

    \theutterance \stepcounter{utterance}  
    & & & \multicolumn{4}{p{0.6\linewidth}}{
        \cellcolor[rgb]{0.9,0.9,0.9}{
            \makecell[{{p{\linewidth}}}]{
                \texttt{\tiny{[GM$\langle$P2]}}
                \texttt{VORSCHLAG: \{'C76', 'C08', 'A03', 'A07', 'C00', 'B38', 'C38', 'C62', 'C10', 'C98'\}} \\
\texttt{ARGUMENT: \{"Durch den Austausch von C56 und C81 gegen C76 und C08 erhöht sich der Gesamtwert deutlich, während der Aufwand noch unter dem Maximum bleibt."\}} \\
            }
        }
    }
     \\ \\

    \theutterance \stepcounter{utterance}  
    & & \multicolumn{4}{p{0.6\linewidth}}{
        \cellcolor[rgb]{0.9,0.9,0.9}{
            \makecell[{{p{\linewidth}}}]{
                \texttt{\tiny{[P1$\langle$GM]}}
                \texttt{VORSCHLAG: \{'C76', 'C08', 'A03', 'A07', 'C00', 'B38', 'C38', 'C62', 'C10', 'C98'\}} \\
\texttt{ARGUMENT: \{"Durch den Austausch von C56 und C81 gegen C76 und C08 erhöht sich der Gesamtwert deutlich, während der Aufwand noch unter dem Maximum bleibt."\}} \\
            }
        }
    }
    & \\ \\

    \theutterance \stepcounter{utterance}  
    & \multicolumn{4}{p{0.6\linewidth}}{
        \cellcolor[rgb]{0.9,0.9,0.9}{
            \makecell[{{p{\linewidth}}}]{
                \texttt{\tiny{[P1$\rangle$GM]}}
                \texttt{ABLEHNUNG: \{'C76', 'C08', 'A03', 'A07', 'C00', 'B38', 'C38', 'C62', 'C10', 'C98'\}} \\
\texttt{ARGUMENT: \{"Das Austauschen von C56 und C81 gegen C76 und C08 reduziert meinen Nutzen erheblich, da C56 und C81 für mich wesentlich wichtiger sind. Ich schlage vor, stattdessen C56 und C81 wieder einzubeziehen und auf andere Gegenstände mit hohem Aufwand zu verzichten."\}} \\
            }
        }
    }
    & & \\ \\

    \theutterance \stepcounter{utterance}  
    & & \multicolumn{4}{p{0.6\linewidth}}{
        \cellcolor[rgb]{0.9,0.9,0.9}{
            \makecell[{{p{\linewidth}}}]{
                \texttt{\tiny{[GM$\rangle$P2]}}
                \texttt{ABLEHNUNG: \{'C76', 'C08', 'A03', 'A07', 'C00', 'B38', 'C38', 'C62', 'C10', 'C98'\}} \\
\texttt{ARGUMENT: \{"Das Austauschen von C56 und C81 gegen C76 und C08 reduziert meinen Nutzen erheblich, da C56 und C81 für mich wesentlich wichtiger sind. Ich schlage vor, stattdessen C56 und C81 wieder einzubeziehen und auf andere Gegenstände mit hohem Aufwand zu verzichten."\}} \\
            }
        }
    }
    & \\ \\

    \theutterance \stepcounter{utterance}  
    & & & \multicolumn{4}{p{0.6\linewidth}}{
        \cellcolor[rgb]{0.9,0.9,0.9}{
            \makecell[{{p{\linewidth}}}]{
                \texttt{\tiny{[GM$\langle$P2]}}
                \texttt{VORSCHLAG: \{'C56', 'C81', 'C08', 'C76', 'C98', 'C10', 'B38', 'C62'\}} \\
\texttt{ARGUMENT: \{"C56 und C81 sind nun enthalten. Dieser Vorschlag maximiert meinen Gesamtwert bei Einhaltung des Maximalaufwands."\}} \\
            }
        }
    }
     \\ \\

    \theutterance \stepcounter{utterance}  
    & & \multicolumn{4}{p{0.6\linewidth}}{
        \cellcolor[rgb]{0.9,0.9,0.9}{
            \makecell[{{p{\linewidth}}}]{
                \texttt{\tiny{[P1$\langle$GM]}}
                \texttt{VORSCHLAG: \{'C56', 'C81', 'C08', 'C76', 'C98', 'C10', 'B38', 'C62'\}} \\
\texttt{ARGUMENT: \{"C56 und C81 sind nun enthalten. Dieser Vorschlag maximiert meinen Gesamtwert bei Einhaltung des Maximalaufwands."\}} \\
            }
        }
    }
    & \\ \\

    \theutterance \stepcounter{utterance}  
    & \multicolumn{4}{p{0.6\linewidth}}{
        \cellcolor[rgb]{0.9,0.9,0.9}{
            \makecell[{{p{\linewidth}}}]{
                \texttt{\tiny{[P1$\rangle$GM]}}
                \texttt{ABLEHNUNG: \{'C56', 'C81', 'C08', 'C76', 'C98', 'C10', 'B38', 'C62'\}} \\
\texttt{ARGUMENT: \{"C08 hat einen hohen Aufwand von 716 und eine geringe Wertigkeit von 121. Bessere Optionen wie A03, A07 und C00 wurden ausgeschlossen."\}} \\
\texttt{VORSCHLAG: \{'A03', 'A07', 'C00', 'C56', 'C38', 'C32', 'C10', 'B38', 'C62', 'C98'\}} \\
            }
        }
    }
    & & \\ \\

    \theutterance \stepcounter{utterance}  
    & & \multicolumn{4}{p{0.6\linewidth}}{
        \cellcolor[rgb]{0.9,0.9,0.9}{
            \makecell[{{p{\linewidth}}}]{
                \texttt{\tiny{[GM$\rangle$P2]}}
                \texttt{ABLEHNUNG: \{'C56', 'C81', 'C08', 'C76', 'C98', 'C10', 'B38', 'C62'\}} \\
\texttt{ARGUMENT: \{"C08 hat einen hohen Aufwand von 716 und eine geringe Wertigkeit von 121. Bessere Optionen wie A03, A07 und C00 wurden ausgeschlossen."\}} \\
\texttt{VORSCHLAG: \{'A03', 'A07', 'C00', 'C56', 'C38', 'C32', 'C10', 'B38', 'C62', 'C98'\}} \\
            }
        }
    }
    & \\ \\

    \theutterance \stepcounter{utterance}  
    & & & \multicolumn{4}{p{0.6\linewidth}}{
        \cellcolor[rgb]{0.9,0.9,0.9}{
            \makecell[{{p{\linewidth}}}]{
                \texttt{\tiny{[GM$\langle$P2]}}
                \texttt{ABLEHNUNG: \{'A03', 'A07', 'C00', 'C56', 'C38', 'C32', 'C10', 'B38', 'C62', 'C98'\}} \\
\texttt{ARGUMENT: \{"C32 hat einen hohen Aufwand von 720 und eine geringe Wertigkeit von 138. Zudem fehlen wichtige Items wie A83 (868 Wert) und C08 (822 Wert), die einen höheren Nutzen bieten."\}} \\
\texttt{VORSCHLAG: \{'C10', 'C98', 'A07', 'C62', 'C00', 'C76', 'C08', 'A83'\}} \\
            }
        }
    }
     \\ \\

    \theutterance \stepcounter{utterance}  
    & & \multicolumn{4}{p{0.6\linewidth}}{
        \cellcolor[rgb]{0.9,0.9,0.9}{
            \makecell[{{p{\linewidth}}}]{
                \texttt{\tiny{[P1$\langle$GM]}}
                \texttt{ABLEHNUNG: \{'A03', 'A07', 'C00', 'C56', 'C38', 'C32', 'C10', 'B38', 'C62', 'C98'\}} \\
\texttt{ARGUMENT: \{"C32 hat einen hohen Aufwand von 720 und eine geringe Wertigkeit von 138. Zudem fehlen wichtige Items wie A83 (868 Wert) und C08 (822 Wert), die einen höheren Nutzen bieten."\}} \\
\texttt{VORSCHLAG: \{'C10', 'C98', 'A07', 'C62', 'C00', 'C76', 'C08', 'A83'\}} \\
            }
        }
    }
    & \\ \\

    \theutterance \stepcounter{utterance}  
    & \multicolumn{4}{p{0.6\linewidth}}{
        \cellcolor[rgb]{0.9,0.9,0.9}{
            \makecell[{{p{\linewidth}}}]{
                \texttt{\tiny{[P1$\rangle$GM]}}
                \texttt{ABLEHNUNG: \{'C10', 'C98', 'A07', 'C62', 'C00', 'C76', 'C08', 'A83'\}} \\
\texttt{ARGUMENT: \{"C08 (Aufwand 716, Wert 121) und A83 (Aufwand 909, Wert 65) haben einen zu hohen Aufwand im Verhältnis zu ihrem Wert. Stattdessen sollten Items wie A03 (Wert 822) und C56 (Wert 868) einbezogen werden, die für mich deutlich höheren Nutzen bieten."\}} \\
\texttt{VORSCHLAG: \{'A03', 'C56', 'A07', 'C00', 'B38', 'C38', 'C62', 'C32', 'C98', 'C10'\}} \\
            }
        }
    }
    & & \\ \\

    \theutterance \stepcounter{utterance}  
    & & \multicolumn{4}{p{0.6\linewidth}}{
        \cellcolor[rgb]{0.9,0.9,0.9}{
            \makecell[{{p{\linewidth}}}]{
                \texttt{\tiny{[GM$\rangle$P2]}}
                \texttt{ABLEHNUNG: \{'C10', 'C98', 'A07', 'C62', 'C00', 'C76', 'C08', 'A83'\}} \\
\texttt{ARGUMENT: \{"C08 (Aufwand 716, Wert 121) und A83 (Aufwand 909, Wert 65) haben einen zu hohen Aufwand im Verhältnis zu ihrem Wert. Stattdessen sollten Items wie A03 (Wert 822) und C56 (Wert 868) einbezogen werden, die für mich deutlich höheren Nutzen bieten."\}} \\
\texttt{VORSCHLAG: \{'A03', 'C56', 'A07', 'C00', 'B38', 'C38', 'C62', 'C32', 'C98', 'C10'\}} \\
            }
        }
    }
    & \\ \\

    \theutterance \stepcounter{utterance}  
    & & & \multicolumn{4}{p{0.6\linewidth}}{
        \cellcolor[rgb]{0.9,0.9,0.9}{
            \makecell[{{p{\linewidth}}}]{
                \texttt{\tiny{[GM$\langle$P2]}}
                \texttt{ABLEHNUNG: \{'A03', 'C56', 'A07', 'C00', 'B38', 'C38', 'C62', 'C32', 'C98', 'C10'\}} \\
\texttt{ARGUMENT: \{"C32 hat einen hohen Aufwand von 720 und für mich nur eine geringe Wichtigkeit. Zudem fehlen Items wie A83 und C08, die für mich deutlich höhere Priorität haben."\}} \\
\texttt{VORSCHLAG: \{'A83', 'C08', 'C76', 'C98', 'C03', 'C10', 'A07'\}} \\
\texttt{ARGUMENT: \{"Diese Kombination priorisiert hochwichtige Gegenstände wie A83, C08 und C76, während Aufwandsintensive Items mit geringer persönlicher Relevanz vermieden wurden. Der Gesamtaufwand liegt unter dem Limit."\}} \\
            }
        }
    }
     \\ \\

    \theutterance \stepcounter{utterance}  
    & & \multicolumn{4}{p{0.6\linewidth}}{
        \cellcolor[rgb]{0.9,0.9,0.9}{
            \makecell[{{p{\linewidth}}}]{
                \texttt{\tiny{[P1$\langle$GM]}}
                \texttt{ABLEHNUNG: \{'A03', 'C56', 'A07', 'C00', 'B38', 'C38', 'C62', 'C32', 'C98', 'C10'\}} \\
\texttt{ARGUMENT: \{"C32 hat einen hohen Aufwand von 720 und für mich nur eine geringe Wichtigkeit. Zudem fehlen Items wie A83 und C08, die für mich deutlich höhere Priorität haben."\}} \\
\texttt{VORSCHLAG: \{'A83', 'C08', 'C76', 'C98', 'C03', 'C10', 'A07'\}} \\
\texttt{ARGUMENT: \{"Diese Kombination priorisiert hochwichtige Gegenstände wie A83, C08 und C76, während Aufwandsintensive Items mit geringer persönlicher Relevanz vermieden wurden. Der Gesamtaufwand liegt unter dem Limit."\}} \\
            }
        }
    }
    & \\ \\

    \theutterance \stepcounter{utterance}  
    & \multicolumn{4}{p{0.6\linewidth}}{
        \cellcolor[rgb]{0.9,0.9,0.9}{
            \makecell[{{p{\linewidth}}}]{
                \texttt{\tiny{[P1$\rangle$GM]}}
                \texttt{ABLEHNUNG: \{'A83', 'C08', 'C76', 'C98', 'C03', 'C10', 'A07'\}} \\
\texttt{ARGUMENT: \{"Die enthaltenen Gegenstände A83, C08 und C76 haben für mich eine sehr niedrige Wertigkeit. Zusätzlich fehlen hochwertige Gegenstände wie C32, C56 und A03, die einen deutlich höheren Nutzen bieten. Der vorgeschlagene Deal ist für mich nicht akzeptabel."\}} \\
\texttt{VORSCHLAG: \{'C32', 'C56', 'A03', 'A07', 'C00', 'B38', 'C38', 'C62', 'C10', 'C98'\}} \\
\texttt{ARGUMENT: \{"Diese Kombination enthält alle für mich besonders wichtigen Gegenstände und bleibt unter dem maximalen Aufwand."\}} \\
            }
        }
    }
    & & \\ \\

    \theutterance \stepcounter{utterance}  
    & & \multicolumn{4}{p{0.6\linewidth}}{
        \cellcolor[rgb]{0.9,0.9,0.9}{
            \makecell[{{p{\linewidth}}}]{
                \texttt{\tiny{[GM$\rangle$P2]}}
                \texttt{ABLEHNUNG: \{'A83', 'C08', 'C76', 'C98', 'C03', 'C10', 'A07'\}} \\
\texttt{ARGUMENT: \{"Die enthaltenen Gegenstände A83, C08 und C76 haben für mich eine sehr niedrige Wertigkeit. Zusätzlich fehlen hochwertige Gegenstände wie C32, C56 und A03, die einen deutlich höheren Nutzen bieten. Der vorgeschlagene Deal ist für mich nicht akzeptabel."\}} \\
\texttt{VORSCHLAG: \{'C32', 'C56', 'A03', 'A07', 'C00', 'B38', 'C38', 'C62', 'C10', 'C98'\}} \\
\texttt{ARGUMENT: \{"Diese Kombination enthält alle für mich besonders wichtigen Gegenstände und bleibt unter dem maximalen Aufwand."\}} \\
            }
        }
    }
    & \\ \\

    \theutterance \stepcounter{utterance}  
    & & & \multicolumn{4}{p{0.6\linewidth}}{
        \cellcolor[rgb]{0.9,0.9,0.9}{
            \makecell[{{p{\linewidth}}}]{
                \texttt{\tiny{[GM$\langle$P2]}}
                \texttt{ABLEHNUNG: \{'C32', 'C56', 'A03', 'A07', 'C00', 'B38', 'C38', 'C62', 'C10', 'C98'\}} \\
\texttt{ARGUMENT: \{"C32 hat einen hohen Aufwand von 720 für nur 138 Wert. Items wie A83 (868) und C08 (822) mit hoher Wertigkeit wurden ausgeschlossen, während Items wie C56 mit nur 65 Wert trotzdem enthalten sind. Bitte überprüfen Sie Ihre Priorisierung."\}} \\
\texttt{VORSCHLAG: \{'A83', 'C08', 'C76', 'C98', 'C62', 'C38', 'C10', 'C00'\}} \\
            }
        }
    }
     \\ \\

    \theutterance \stepcounter{utterance}  
    & & \multicolumn{4}{p{0.6\linewidth}}{
        \cellcolor[rgb]{0.9,0.9,0.9}{
            \makecell[{{p{\linewidth}}}]{
                \texttt{\tiny{[P1$\langle$GM]}}
                \texttt{ABLEHNUNG: \{'C32', 'C56', 'A03', 'A07', 'C00', 'B38', 'C38', 'C62', 'C10', 'C98'\}} \\
\texttt{ARGUMENT: \{"C32 hat einen hohen Aufwand von 720 für nur 138 Wert. Items wie A83 (868) und C08 (822) mit hoher Wertigkeit wurden ausgeschlossen, während Items wie C56 mit nur 65 Wert trotzdem enthalten sind. Bitte überprüfen Sie Ihre Priorisierung."\}} \\
\texttt{VORSCHLAG: \{'A83', 'C08', 'C76', 'C98', 'C62', 'C38', 'C10', 'C00'\}} \\
            }
        }
    }
    & \\ \\

    \theutterance \stepcounter{utterance}  
    & \multicolumn{4}{p{0.6\linewidth}}{
        \cellcolor[rgb]{0.9,0.9,0.9}{
            \makecell[{{p{\linewidth}}}]{
                \texttt{\tiny{[P1$\rangle$GM]}}
                \texttt{ABLEHNUNG: \{'A83', 'C08', 'C76', 'C98', 'C62', 'C38', 'C10', 'C00'\}} \\
\texttt{ARGUMENT: \{"A83 (65), C08 (121) und C76 (138) haben extrem geringe Wichtigkeit für mich. Die fehlenden Items A03 (822), C56 (868) und C32 (808) bringen erheblich mehr Nutzen."\}} \\
\texttt{VORSCHLAG: \{'A03', 'C56', 'C32', 'A07', 'C00', 'B38', 'C38', 'C62', 'C10', 'C98'\}} \\
            }
        }
    }
    & & \\ \\

    \theutterance \stepcounter{utterance}  
    & & \multicolumn{4}{p{0.6\linewidth}}{
        \cellcolor[rgb]{0.9,0.9,0.9}{
            \makecell[{{p{\linewidth}}}]{
                \texttt{\tiny{[GM$\rangle$P2]}}
                \texttt{ABLEHNUNG: \{'A83', 'C08', 'C76', 'C98', 'C62', 'C38', 'C10', 'C00'\}} \\
\texttt{ARGUMENT: \{"A83 (65), C08 (121) und C76 (138) haben extrem geringe Wichtigkeit für mich. Die fehlenden Items A03 (822), C56 (868) und C32 (808) bringen erheblich mehr Nutzen."\}} \\
\texttt{VORSCHLAG: \{'A03', 'C56', 'C32', 'A07', 'C00', 'B38', 'C38', 'C62', 'C10', 'C98'\}} \\
            }
        }
    }
    & \\ \\

    \theutterance \stepcounter{utterance}  
    & & & \multicolumn{4}{p{0.6\linewidth}}{
        \cellcolor[rgb]{0.9,0.9,0.9}{
            \makecell[{{p{\linewidth}}}]{
                \texttt{\tiny{[GM$\langle$P2]}}
                \texttt{ABLEHNUNG: \{'A03', 'C56', 'C32', 'A07', 'C00', 'B38', 'C38', 'C62', 'C10', 'C98'\}} \\
\texttt{ARGUMENT: \{"C32 hat einen hohen Aufwand von 720 und für mich nur eine Wertigkeit von 138. Gegenstände wie A83 (868 Wert) und C08 (822 Wert) haben eine deutlich höhere Priorität für mich und sollten anstelle von C32 bevorzugt werden."\}} \\
            }
        }
    }
     \\ \\

    \theutterance \stepcounter{utterance}  
    & & \multicolumn{4}{p{0.6\linewidth}}{
        \cellcolor[rgb]{0.9,0.9,0.9}{
            \makecell[{{p{\linewidth}}}]{
                \texttt{\tiny{[P1$\langle$GM]}}
                \texttt{ABLEHNUNG: \{'A03', 'C56', 'C32', 'A07', 'C00', 'B38', 'C38', 'C62', 'C10', 'C98'\}} \\
\texttt{ARGUMENT: \{"C32 hat einen hohen Aufwand von 720 und für mich nur eine Wertigkeit von 138. Gegenstände wie A83 (868 Wert) und C08 (822 Wert) haben eine deutlich höhere Priorität für mich und sollten anstelle von C32 bevorzugt werden."\}} \\
            }
        }
    }
    & \\ \\

    \theutterance \stepcounter{utterance}  
    & \multicolumn{4}{p{0.6\linewidth}}{
        \cellcolor[rgb]{0.9,0.9,0.9}{
            \makecell[{{p{\linewidth}}}]{
                \texttt{\tiny{[P1$\rangle$GM]}}
                \texttt{VORSCHLAG: \{'A03', 'C56', 'A07', 'C00', 'B38', 'C38', 'C62', 'C10', 'C98', 'C08'\}} \\
\texttt{ARGUMENT: \{"C08 wurde mit hoher Priorität für Sie aufgenommen, während C32 entfernt wurde. Dieser Vorschlag bietet eine faire Aufteilung, die unter dem Maximalaufwand bleibt und beide Parteien zufriedenstellen kann."\}} \\
            }
        }
    }
    & & \\ \\

    \theutterance \stepcounter{utterance}  
    & & \multicolumn{4}{p{0.6\linewidth}}{
        \cellcolor[rgb]{0.9,0.9,0.9}{
            \makecell[{{p{\linewidth}}}]{
                \texttt{\tiny{[GM$\rangle$P2]}}
                \texttt{VORSCHLAG: \{'A03', 'C56', 'A07', 'C00', 'B38', 'C38', 'C62', 'C10', 'C98', 'C08'\}} \\
\texttt{ARGUMENT: \{"C08 wurde mit hoher Priorität für Sie aufgenommen, während C32 entfernt wurde. Dieser Vorschlag bietet eine faire Aufteilung, die unter dem Maximalaufwand bleibt und beide Parteien zufriedenstellen kann."\}} \\
            }
        }
    }
    & \\ \\

    \theutterance \stepcounter{utterance}  
    & & & \multicolumn{4}{p{0.6\linewidth}}{
        \cellcolor[rgb]{0.9,0.9,0.9}{
            \makecell[{{p{\linewidth}}}]{
                \texttt{\tiny{[GM$\langle$P2]}}
                \texttt{ABLEHNUNG: \{'A03', 'C56', 'A07', 'C00', 'B38', 'C38', 'C62', 'C10', 'C98', 'C08'\}} \\
\texttt{ARGUMENT: \{"C56 und A03 haben für mich sehr geringe Priorität, während A83 und C76, die meine höchste Wertigkeit haben, fehlen. Der Vorschlag ignoriert klar meine dringendsten Bedürfnisse."\}} \\
            }
        }
    }
     \\ \\

    \theutterance \stepcounter{utterance}  
    & & \multicolumn{4}{p{0.6\linewidth}}{
        \cellcolor[rgb]{0.9,0.9,0.9}{
            \makecell[{{p{\linewidth}}}]{
                \texttt{\tiny{[P1$\langle$GM]}}
                \texttt{ABLEHNUNG: \{'A03', 'C56', 'A07', 'C00', 'B38', 'C38', 'C62', 'C10', 'C98', 'C08'\}} \\
\texttt{ARGUMENT: \{"C56 und A03 haben für mich sehr geringe Priorität, während A83 und C76, die meine höchste Wertigkeit haben, fehlen. Der Vorschlag ignoriert klar meine dringendsten Bedürfnisse."\}} \\
            }
        }
    }
    & \\ \\

    \theutterance \stepcounter{utterance}  
    & \multicolumn{4}{p{0.6\linewidth}}{
        \cellcolor[rgb]{0.9,0.9,0.9}{
            \makecell[{{p{\linewidth}}}]{
                \texttt{\tiny{[P1$\rangle$GM]}}
                \texttt{VORSCHLAG: \{'A83', 'C76', 'C08', 'C00', 'A07', 'B38', 'C62'\}} \\
\texttt{ARGUMENT: \{"Die Vorschlag beinhaltet Ihre prioritären Gegenstände A83 und C76, kombiniert mit hochwichtigen Items für mich wie C00 und A07. Der Gesamtaufwand liegt deutlich unter dem Limit und bietet eine faire Ausgewogenheit für beide Seiten."\}} \\
            }
        }
    }
    & & \\ \\

    \theutterance \stepcounter{utterance}  
    & & \multicolumn{4}{p{0.6\linewidth}}{
        \cellcolor[rgb]{0.9,0.9,0.9}{
            \makecell[{{p{\linewidth}}}]{
                \texttt{\tiny{[GM$\rangle$P2]}}
                \texttt{VORSCHLAG: \{'A83', 'C76', 'C08', 'C00', 'A07', 'B38', 'C62'\}} \\
\texttt{ARGUMENT: \{"Die Vorschlag beinhaltet Ihre prioritären Gegenstände A83 und C76, kombiniert mit hochwichtigen Items für mich wie C00 und A07. Der Gesamtaufwand liegt deutlich unter dem Limit und bietet eine faire Ausgewogenheit für beide Seiten."\}} \\
            }
        }
    }
    & \\ \\

    \theutterance \stepcounter{utterance}  
    & & & \multicolumn{4}{p{0.6\linewidth}}{
        \cellcolor[rgb]{0.9,0.9,0.9}{
            \makecell[{{p{\linewidth}}}]{
                \texttt{\tiny{[GM$\langle$P2]}}
                \texttt{ABLEHNUNG: \{'A83', 'C76', 'C08', 'C00', 'A07', 'B38', 'C62'\}} \\
\texttt{ARGUMENT: \{"B38 und C00 haben für mich eine niedrige Priorität. Bitte ersetzen Sie sie durch C98 und C81, um den Gesamtwert zu erhöhen."\}} \\
            }
        }
    }
     \\ \\

    \theutterance \stepcounter{utterance}  
    & & \multicolumn{4}{p{0.6\linewidth}}{
        \cellcolor[rgb]{0.9,0.9,0.9}{
            \makecell[{{p{\linewidth}}}]{
                \texttt{\tiny{[P1$\langle$GM]}}
                \texttt{ABLEHNUNG: \{'A83', 'C76', 'C08', 'C00', 'A07', 'B38', 'C62'\}} \\
\texttt{ARGUMENT: \{"B38 und C00 haben für mich eine niedrige Priorität. Bitte ersetzen Sie sie durch C98 und C81, um den Gesamtwert zu erhöhen."\}} \\
            }
        }
    }
    & \\ \\

    \theutterance \stepcounter{utterance}  
    & \multicolumn{4}{p{0.6\linewidth}}{
        \cellcolor[rgb]{0.9,0.9,0.9}{
            \makecell[{{p{\linewidth}}}]{
                \texttt{\tiny{[P1$\rangle$GM]}}
                \texttt{VORSCHLAG: \{'A83', 'C76', 'C56', 'A07', 'C00', 'C62', 'C38', 'C98', 'C10'\}} \\
\texttt{ARGUMENT: \{"Diese Kombination enthält Ihre prioritären Gegenstände A83 und C76 sowie meine Top{-}Prioritäten C56, A07 und C00. Der Gesamtaufwand beträgt 3105 und ist deutlich unter dem Limit, was eine faire Ausgewogenheit für beide Seiten darstellt."\}} \\
            }
        }
    }
    & & \\ \\

    \theutterance \stepcounter{utterance}  
    & & \multicolumn{4}{p{0.6\linewidth}}{
        \cellcolor[rgb]{0.9,0.9,0.9}{
            \makecell[{{p{\linewidth}}}]{
                \texttt{\tiny{[GM$\rangle$P2]}}
                \texttt{VORSCHLAG: \{'A83', 'C76', 'C56', 'A07', 'C00', 'C62', 'C38', 'C98', 'C10'\}} \\
\texttt{ARGUMENT: \{"Diese Kombination enthält Ihre prioritären Gegenstände A83 und C76 sowie meine Top{-}Prioritäten C56, A07 und C00. Der Gesamtaufwand beträgt 3105 und ist deutlich unter dem Limit, was eine faire Ausgewogenheit für beide Seiten darstellt."\}} \\
            }
        }
    }
    & \\ \\

    \theutterance \stepcounter{utterance}  
    & & & \multicolumn{4}{p{0.6\linewidth}}{
        \cellcolor[rgb]{0.9,0.9,0.9}{
            \makecell[{{p{\linewidth}}}]{
                \texttt{\tiny{[GM$\langle$P2]}}
                \texttt{ABLEHNUNG: \{'A83', 'C76', 'C56', 'A07', 'C00', 'C62', 'C38', 'C98', 'C10'\}} \\
\texttt{ARGUMENT: \{"C56 hat für mich eine sehr niedrige Wertigkeit, während C08 mit einer sehr hohen Priorität fehlt. Bitte ersetze C56 durch C08."\}} \\
            }
        }
    }
     \\ \\

    \theutterance \stepcounter{utterance}  
    & & \multicolumn{4}{p{0.6\linewidth}}{
        \cellcolor[rgb]{0.9,0.9,0.9}{
            \makecell[{{p{\linewidth}}}]{
                \texttt{\tiny{[P1$\langle$GM]}}
                \texttt{ABLEHNUNG: \{'A83', 'C76', 'C56', 'A07', 'C00', 'C62', 'C38', 'C98', 'C10'\}} \\
\texttt{ARGUMENT: \{"C56 hat für mich eine sehr niedrige Wertigkeit, während C08 mit einer sehr hohen Priorität fehlt. Bitte ersetze C56 durch C08."\}} \\
            }
        }
    }
    & \\ \\

    \theutterance \stepcounter{utterance}  
    & \multicolumn{4}{p{0.6\linewidth}}{
        \cellcolor[rgb]{0.9,0.9,0.9}{
            \makecell[{{p{\linewidth}}}]{
                \texttt{\tiny{[P1$\rangle$GM]}}
                \texttt{VORSCHLAG: \{'A83', 'C08', 'C81', 'A03', 'A07', 'C62', 'C10', 'C98'\}} \\
\texttt{ARGUMENT: \{"Wir schlagen A83, C08, C81, A03, A07, C62, C10 und C98 vor. Dies berücksichtigt Ihre Prioritäten für A83 und C08 sowie C98 und C81. Gleichzeitig maximieren wir meine hochwertigen Items wie A03 und A07. Der Gesamtaufwand beträgt 3027 und ist damit unter dem Limit."\}} \\
            }
        }
    }
    & & \\ \\

    \theutterance \stepcounter{utterance}  
    & & \multicolumn{4}{p{0.6\linewidth}}{
        \cellcolor[rgb]{0.9,0.9,0.9}{
            \makecell[{{p{\linewidth}}}]{
                \texttt{\tiny{[GM$\rangle$P2]}}
                \texttt{VORSCHLAG: \{'A83', 'C08', 'C81', 'A03', 'A07', 'C62', 'C10', 'C98'\}} \\
\texttt{ARGUMENT: \{"Wir schlagen A83, C08, C81, A03, A07, C62, C10 und C98 vor. Dies berücksichtigt Ihre Prioritäten für A83 und C08 sowie C98 und C81. Gleichzeitig maximieren wir meine hochwertigen Items wie A03 und A07. Der Gesamtaufwand beträgt 3027 und ist damit unter dem Limit."\}} \\
            }
        }
    }
    & \\ \\

    \theutterance \stepcounter{utterance}  
    & & & \multicolumn{4}{p{0.6\linewidth}}{
        \cellcolor[rgb]{0.9,0.9,0.9}{
            \makecell[{{p{\linewidth}}}]{
                \texttt{\tiny{[GM$\langle$P2]}}
                \texttt{ABLEHNUNG: \{'A83', 'C08', 'C81', 'A03', 'A07', 'C62', 'C10', 'C98'\}} \\
\texttt{ARGUMENT: \{"C81 (583) und A03 (121) haben eine zu niedrige Wertigkeit für mich. Ersetze sie durch C76 (808) und C00 (389) für einen höheren Gesamtwert."\}} \\
            }
        }
    }
     \\ \\

    \theutterance \stepcounter{utterance}  
    & & \multicolumn{4}{p{0.6\linewidth}}{
        \cellcolor[rgb]{0.9,0.9,0.9}{
            \makecell[{{p{\linewidth}}}]{
                \texttt{\tiny{[P1$\langle$GM]}}
                \texttt{ABLEHNUNG: \{'A83', 'C08', 'C81', 'A03', 'A07', 'C62', 'C10', 'C98'\}} \\
\texttt{ARGUMENT: \{"C81 (583) und A03 (121) haben eine zu niedrige Wertigkeit für mich. Ersetze sie durch C76 (808) und C00 (389) für einen höheren Gesamtwert."\}} \\
            }
        }
    }
    & \\ \\

    \theutterance \stepcounter{utterance}  
    & \multicolumn{4}{p{0.6\linewidth}}{
        \cellcolor[rgb]{0.9,0.9,0.9}{
            \makecell[{{p{\linewidth}}}]{
                \texttt{\tiny{[P1$\rangle$GM]}}
                \texttt{VORSCHLAG: \{'A83', 'C08', 'C76', 'C00', 'C56', 'C62', 'C10', 'C98'\}} \\
\texttt{ARGUMENT: \{"Beinhaltet Ihre prioritären Gegenstände A83 und C08 sowie meine Top{-}High{-}Wert{-}Items wie C00 und C56. Der Gesamtaufwand von 3198 liegt deutlich unter dem Limit und garantiert faire Gewinne für beide Parteien."\}} \\
            }
        }
    }
    & & \\ \\

    \theutterance \stepcounter{utterance}  
    & & \multicolumn{4}{p{0.6\linewidth}}{
        \cellcolor[rgb]{0.9,0.9,0.9}{
            \makecell[{{p{\linewidth}}}]{
                \texttt{\tiny{[GM$\rangle$P2]}}
                \texttt{VORSCHLAG: \{'A83', 'C08', 'C76', 'C00', 'C56', 'C62', 'C10', 'C98'\}} \\
\texttt{ARGUMENT: \{"Beinhaltet Ihre prioritären Gegenstände A83 und C08 sowie meine Top{-}High{-}Wert{-}Items wie C00 und C56. Der Gesamtaufwand von 3198 liegt deutlich unter dem Limit und garantiert faire Gewinne für beide Parteien."\}} \\
            }
        }
    }
    & \\ \\

    \theutterance \stepcounter{utterance}  
    & & & \multicolumn{4}{p{0.6\linewidth}}{
        \cellcolor[rgb]{0.9,0.9,0.9}{
            \makecell[{{p{\linewidth}}}]{
                \texttt{\tiny{[GM$\langle$P2]}}
                \texttt{ABLEHNUNG: \{'A83', 'C08', 'C76', 'C00', 'C56', 'C62', 'C10', 'C98'\}} \\
\texttt{ARGUMENT: \{"C56 hat eine sehr niedrige Wertigkeit für mich. Bitte ersetze C56 durch C38, um den Gesamtwert zu erhöhen."\}} \\
            }
        }
    }
     \\ \\

    \theutterance \stepcounter{utterance}  
    & & \multicolumn{4}{p{0.6\linewidth}}{
        \cellcolor[rgb]{0.9,0.9,0.9}{
            \makecell[{{p{\linewidth}}}]{
                \texttt{\tiny{[P1$\langle$GM]}}
                \texttt{ABLEHNUNG: \{'A83', 'C08', 'C76', 'C00', 'C56', 'C62', 'C10', 'C98'\}} \\
\texttt{ARGUMENT: \{"C56 hat eine sehr niedrige Wertigkeit für mich. Bitte ersetze C56 durch C38, um den Gesamtwert zu erhöhen."\}} \\
            }
        }
    }
    & \\ \\

    \theutterance \stepcounter{utterance}  
    & \multicolumn{4}{p{0.6\linewidth}}{
        \cellcolor[rgb]{0.9,0.9,0.9}{
            \makecell[{{p{\linewidth}}}]{
                \texttt{\tiny{[P1$\rangle$GM]}}
                \texttt{VORSCHLAG: \{'C08', 'C76', 'C00', 'C38', 'C62', 'C10', 'C98', 'B38', 'A03'\}} \\
\texttt{ARGUMENT: \{"C56 wurde wie gewünscht entfernt und durch C38 ersetzt. Die neue Kombination umfasst hochwertige Items für beide Seiten wie A03 und B38, bleibt unter dem Maximalaufwand und stellt eine ausgewogene Lösung dar."\}} \\
            }
        }
    }
    & & \\ \\

    \theutterance \stepcounter{utterance}  
    & & \multicolumn{4}{p{0.6\linewidth}}{
        \cellcolor[rgb]{0.9,0.9,0.9}{
            \makecell[{{p{\linewidth}}}]{
                \texttt{\tiny{[GM$\rangle$P2]}}
                \texttt{VORSCHLAG: \{'C08', 'C76', 'C00', 'C38', 'C62', 'C10', 'C98', 'B38', 'A03'\}} \\
\texttt{ARGUMENT: \{"C56 wurde wie gewünscht entfernt und durch C38 ersetzt. Die neue Kombination umfasst hochwertige Items für beide Seiten wie A03 und B38, bleibt unter dem Maximalaufwand und stellt eine ausgewogene Lösung dar."\}} \\
            }
        }
    }
    & \\ \\

    \theutterance \stepcounter{utterance}  
    & & & \multicolumn{4}{p{0.6\linewidth}}{
        \cellcolor[rgb]{0.9,0.9,0.9}{
            \makecell[{{p{\linewidth}}}]{
                \texttt{\tiny{[GM$\langle$P2]}}
                \texttt{ABLEHNUNG: \{'C08', 'C76', 'C00', 'C38', 'C62', 'C10', 'C98', 'B38', 'A03'\}} \\
\texttt{ARGUMENT: \{"A83 mit 868 Wert ist meine höchste Priorität und muss enthalten sein. Bitte ersetze B38 und A03 mit A83, um dies zu ermöglichen."\}} \\
            }
        }
    }
     \\ \\

    \theutterance \stepcounter{utterance}  
    & & \multicolumn{4}{p{0.6\linewidth}}{
        \cellcolor[rgb]{0.9,0.9,0.9}{
            \makecell[{{p{\linewidth}}}]{
                \texttt{\tiny{[P1$\langle$GM]}}
                \texttt{ABLEHNUNG: \{'C08', 'C76', 'C00', 'C38', 'C62', 'C10', 'C98', 'B38', 'A03'\}} \\
\texttt{ARGUMENT: \{"A83 mit 868 Wert ist meine höchste Priorität und muss enthalten sein. Bitte ersetze B38 und A03 mit A83, um dies zu ermöglichen."\}} \\
            }
        }
    }
    & \\ \\

    \theutterance \stepcounter{utterance}  
    & \multicolumn{4}{p{0.6\linewidth}}{
        \cellcolor[rgb]{0.9,0.9,0.9}{
            \makecell[{{p{\linewidth}}}]{
                \texttt{\tiny{[P1$\rangle$GM]}}
                \texttt{VORSCHLAG: \{'A83', 'C08', 'C56', 'C00', 'A03', 'C38', 'C62', 'C10', 'C98'\}} \\
\texttt{ARGUMENT: \{"Dieser Vorschlag enthält Ihre Prioritäten A83 und C08 sowie meine Top{-}Werte wie C56, A03 und C00. Der Gesamtaufwand beträgt 3306 und liegt deutlich unter dem Maximum von 3357."\}} \\
            }
        }
    }
    & & \\ \\

    \theutterance \stepcounter{utterance}  
    & & \multicolumn{4}{p{0.6\linewidth}}{
        \cellcolor[rgb]{0.9,0.9,0.9}{
            \makecell[{{p{\linewidth}}}]{
                \texttt{\tiny{[GM$\rangle$P2]}}
                \texttt{VORSCHLAG: \{'A83', 'C08', 'C56', 'C00', 'A03', 'C38', 'C62', 'C10', 'C98'\}} \\
\texttt{ARGUMENT: \{"Dieser Vorschlag enthält Ihre Prioritäten A83 und C08 sowie meine Top{-}Werte wie C56, A03 und C00. Der Gesamtaufwand beträgt 3306 und liegt deutlich unter dem Maximum von 3357."\}} \\
            }
        }
    }
    & \\ \\

    \theutterance \stepcounter{utterance}  
    & & & \multicolumn{4}{p{0.6\linewidth}}{
        \cellcolor[rgb]{0.9,0.9,0.9}{
            \makecell[{{p{\linewidth}}}]{
                \texttt{\tiny{[GM$\langle$P2]}}
                \texttt{ABLEHNUNG: \{'A83', 'C08', 'C56', 'C00', 'A03', 'C38', 'C62', 'C10', 'C98'\}} \\
\texttt{ARGUMENT: \{"C56 (65) und A03 (121) haben für mich eine sehr niedrige Wertigkeit. Bitte ersetze sie durch C76 (808) und C03 (780), um die Gesamt{-}Wertigkeit signifikant zu erhöhen, während der Aufwand unter 3357 bleibt."\}} \\
            }
        }
    }
     \\ \\

    \theutterance \stepcounter{utterance}  
    & & \multicolumn{4}{p{0.6\linewidth}}{
        \cellcolor[rgb]{0.9,0.9,0.9}{
            \makecell[{{p{\linewidth}}}]{
                \texttt{\tiny{[P1$\langle$GM]}}
                \texttt{ABLEHNUNG: \{'A83', 'C08', 'C56', 'C00', 'A03', 'C38', 'C62', 'C10', 'C98'\}} \\
\texttt{ARGUMENT: \{"C56 (65) und A03 (121) haben für mich eine sehr niedrige Wertigkeit. Bitte ersetze sie durch C76 (808) und C03 (780), um die Gesamt{-}Wertigkeit signifikant zu erhöhen, während der Aufwand unter 3357 bleibt."\}} \\
            }
        }
    }
    & \\ \\

    \theutterance \stepcounter{utterance}  
    & \multicolumn{4}{p{0.6\linewidth}}{
        \cellcolor[rgb]{0.9,0.9,0.9}{
            \makecell[{{p{\linewidth}}}]{
                \texttt{\tiny{[P1$\rangle$GM]}}
                \texttt{VORSCHLAG: \{'A83', 'C08', 'C76', 'C00', 'C38', 'C62', 'C10', 'C98'\}} \\
\texttt{ARGUMENT: \{"Dieser Vorschlag enthält Ihre prioritären Gegenstände A83, C08 und C76 sowie wichtige Items für mich wie C00 (Wert 780) und C38 (Wert 583). Der Gesamtaufwand beträgt 3352 und bleibt unter dem Limit, was eine faire Balance bietet."\}} \\
            }
        }
    }
    & & \\ \\

    \theutterance \stepcounter{utterance}  
    & & \multicolumn{4}{p{0.6\linewidth}}{
        \cellcolor[rgb]{0.9,0.9,0.9}{
            \makecell[{{p{\linewidth}}}]{
                \texttt{\tiny{[GM$\rangle$P2]}}
                \texttt{VORSCHLAG: \{'A83', 'C08', 'C76', 'C00', 'C38', 'C62', 'C10', 'C98'\}} \\
\texttt{ARGUMENT: \{"Dieser Vorschlag enthält Ihre prioritären Gegenstände A83, C08 und C76 sowie wichtige Items für mich wie C00 (Wert 780) und C38 (Wert 583). Der Gesamtaufwand beträgt 3352 und bleibt unter dem Limit, was eine faire Balance bietet."\}} \\
            }
        }
    }
    & \\ \\

    \theutterance \stepcounter{utterance}  
    & & & \multicolumn{4}{p{0.6\linewidth}}{
        \cellcolor[rgb]{0.9,0.9,0.9}{
            \makecell[{{p{\linewidth}}}]{
                \texttt{\tiny{[GM$\langle$P2]}}
                \texttt{assistant} \\
\texttt{ZUSTIMMUNG: \{'A83', 'C08', 'C76', 'C00', 'C38', 'C62', 'C10', 'C98'\}} \\
            }
        }
    }
     \\ \\

    \theutterance \stepcounter{utterance}  
    & & \multicolumn{4}{p{0.6\linewidth}}{
        \cellcolor[rgb]{0.9,0.9,0.9}{
            \makecell[{{p{\linewidth}}}]{
                \texttt{\tiny{[GM$\rangle$P2]}}
                \texttt{Deine Antwort hat kein argument enthalten. Du muss das strukturierte Format ARGUMENT: \{...\} irgendwo in theiner Antwort verwenden. Versuche es erneut} \\
            }
        }
    }
    & \\ \\

    \theutterance \stepcounter{utterance}  
    & & & \multicolumn{4}{p{0.6\linewidth}}{
        \cellcolor[rgb]{0.9,0.9,0.9}{
            \makecell[{{p{\linewidth}}}]{
                \texttt{\tiny{[GM$\langle$P2]}}
                \texttt{ZUSTIMMUNG: \{'A83', 'C08', 'C76', 'C00', 'C38', 'C62', 'C10', 'C98'\}} \\
\texttt{ARGUMENT: \{"Der Gesamtaufwand ist unter dem Limit und die ausgewählten Gegenstände decken die wichtigsten Prioritäten beider Seiten ab. Dieser Vorschlag ist für mich akzeptabel."\}} \\
            }
        }
    }
     \\ \\

    \theutterance \stepcounter{utterance}  
    & & & \multicolumn{2}{p{0.3\linewidth}}{
        \cellcolor[rgb]{0.9,0.9,0.9}{
            \makecell[{{p{\linewidth}}}]{
                \texttt{\tiny{[GM$|$GM]}}
                \texttt{game successful} \\
            }
        }
    }
    & & \\ \\

    \theutterance \stepcounter{utterance}  
    & & & \multicolumn{2}{p{0.3\linewidth}}{
        \cellcolor[rgb]{0.9,0.9,0.9}{
            \makecell[{{p{\linewidth}}}]{
                \texttt{\tiny{[GM$|$GM]}}
                \texttt{end game} \\
            }
        }
    }
    & & \\ \\

\end{supertabular}
}

\end{document}
